\begin{rubric}{Research \& Lab Experience}
\entry*[2024 -- Present]%
	\textbf{Sun Yadong Lab, ShanghaiTech University}\par
	Undergraduate Researcher\par
	Systematic training and practice in the \href{https://slst.shanghaitech.edu.cn/syd/main.htm}{Sun Yadong Lab}, which focuses on the structure and function of protein-RNA complexes, primarily using cryo-electron microscopy (cryo-EM) and other structural biology techniques to resolve three-dimensional structures of biological macromolecules. Acquired key skills including molecular cloning, plasmid construction, cell culture, protein expression and purification, as well as purification methods such as affinity chromatography and size exclusion chromatography (SEC). Through participation in daily research activities, gained hands-on experience in structural biology experiments, including sample preparation, data collection and processing, and developed a deeper understanding of the mechanisms underlying protein-RNA interactions.\par
%
\entry*[Fall 2025]%
	\textbf{FlyCropper -- Drosophila Wing Classification System}\par
	Genetics Lab Independent Project\par
	One of the primary code contributors for an independent experimental project in the Genetics laboratory course. The system uses a two-stage deep-learning pipeline (object detection + classification) to automatically detect Drosophila and classify wing phenotypes (long / short wing).\par
	FlyCropper is a deep learning-based Drosophila detection and wing classification system using a two-stage pipeline: YOLOv8 for object detection followed by ResNet/EfficientNet for classification, enabling automated long/short wing identification. Supports the complete workflow from Label Studio annotation to end-to-end inference.\par
	\href{https://github.com/xiangyuw1/FlyCropper}{github.com/xiangyuw1/FlyCropper}
%
\entry*[Spring 2025]%
	\textbf{Evaluating Inhibitory Effects of mTORC1 Inhibitors on Cancer Cell Lines}\par
	Cell Biology Lab Independent Project\par
	Responsible for the design and execution of the biochemistry experimental component of this independent project in the Cell Biology laboratory course, validating the inhibitory effects of mTORC1 inhibitor WRX606 on multiple cancer cell lines by detecting transcriptional changes of downstream target genes FASN and PDK1 via qPCR (with ACTB as reference gene).
%
\entry*[Fall 2025]%
	\textbf{Computational Simulation of Circadian Rhythm Regulation in Cyanobacteria}\par
	Introduction to Systems Biology Course Project\par
	Built an ordinary differential equation (ODE) model of the KaiABC protein phosphorylation dynamics based on the cyanobacterial circadian oscillator system, simulating the core mechanism of the biological clock. The model targets the ordered phosphorylation/dephosphorylation cycle of KaiC at serine (S431) and threonine (T432) residues, regulated by KaiA (activator) and KaiB (inhibitor), producing an autonomous oscillation of approximately 24 hours. Extracted 8 sets of rate constants and KaiA-KaiC binding kinetics parameters from literature, established a four-state (U/T/S/ST-KaiC) ODE system, and completed parameter calculation, numerical simulation, and result visualization using Python (NumPy, SciPy, Matplotlib), successfully reproducing the phosphorylation oscillation curves from the literature.\par
	\href{https://github.com/xiangyuw1/kaic-oscillator}{github.com/xiangyuw1/kaic-oscillator}
%
\end{rubric}
